% sec-02: the five feasibility questions and the single approval step.
% Table 11 and the actionbox.  No figure: the agenda is a table, and drawing it
% would restate it.
\section{The Five Questions a First Meeting Settles}

Every question below is answerable in a 45-minute meeting by people who are
already in the room. None requires a document read in advance, a committee
convened, or a decision escalated. A feasibility meeting whose agenda requires
preparation is a feasibility meeting that gets rescheduled.

\begin{apptable}
\begin{tabularx}{\textwidth}{>{\raggedright\arraybackslash}p{0.5cm} >{\raggedright\arraybackslash}p{4.8cm} >{\raggedright\arraybackslash}X}
\headrow \thead{\#} & \thead{The question} & \thead{What a yes unlocks, and what a no costs}\\
\midrule
1 & Is the population and the perioperative concept clinically supportable here & Everything else on the list. A no means stop rather than negotiate, and stopping early is the cheap outcome \\
2 & Could a pancreatic surgeon here serve as site principal investigator, and a GI oncologist lead systemic therapy & The regulatory structure. The sponsor cannot supply this role, because every 21 CFR 54.2 trigger would fire if one person held both \\
3 & What is the annual eligible-patient count, and what competes for it & An enrollment period a funder can believe. Three competing protocols on the same population turn an 18-participant escalation into a five-year trial \\
4 & Is the robotic platform available, and is the workflow permissible on it & A configuration that can be named in a regulatory submission. Without one, the submission cannot be written at all \\
5 & Would the institution host an investigator-initiated trial sponsored by a small company & The contracting, budgeting and review workstreams. It is asked last because it is most likely to be answered by policy rather than by the room \\
\end{tabularx}
\end{apptable}
\vspace{-0.60cm}
\tabcap{Table 11. The five feasibility questions in the order a meeting should\\
take them, with what a yes to each unlocks and what a no to each costs.}

Question 4 is where the robotic configuration gets specified. No manufacturer,
model, arm count, or software version is asserted anywhere by this company: the
configuration is settled at a site agreement, and this question is the
conversation that settles it. What is needed from the site is which systems are
available, which instruments, which software versions, what networking
restrictions apply, what cybersecurity requirements apply, and whether a recorded
research session is permissible.

\subsection{The one item that gates all five}

Which party is capable of serving as sponsor-investigator. The applicant, the
site principal investigator, the FDA sponsor, the investigational new drug
application holder, and the device manufacturer do not have to be the same
entity, but their responsibilities have to be explicit, and today they are
consistent only with this company's assumption about them
\cite{ecfr2026part312,ucsdactri2026}.

That determination is requested separately from clinical trial support services,
because it is a regulatory question rather than a clinical one and it should not
consume a clinical meeting.

\subsection{What is brought to the meeting, and what is not asked for}

A Phase 1 protocol, an investigational new drug application, a Phase 2 protocol,
quantitative simulations with their verification notebooks, and a fifteen-document
site package are complete, public, and deposited
\cite{onpremwhippl,phase1ind,phase2proto,qspmetpancre,movein2026}. The budget
frame is \$700,000 per year and \$3,500,000 over five years in direct cost
\cite{tenapps2026}, and \$420,000 of private capital is allocated to site
start-up and committed to nobody.

What is not asked for in a first meeting: an endorsement, a hosting commitment, a
letter of support for a grant application, a decision from senior leadership
before a surgeon and an oncologist have formed a view, or the sharing of any
confidential engineering detail before a confidentiality agreement is in place.

\subsection{What happens in the week after the meeting, by outcome}

A feasibility meeting has three possible outcomes and each has a different next
move, written now so that none of them costs a week of deliberation later.

If the answer to questions 1, 2 and 5 is yes, the concept intake pack in
\appfile{funding/auto-fund/04Sep26/forms} is submitted with the named
investigator, the budget questionnaire is returned to the trials office, and
stage 1 of the site start-up draw in \S5 releases against the written response.

If the answer to question 1 or 2 is no, the right move is to stop with that
institution rather than to negotiate. A population that a pancreatic surgical
oncologist judges unsupportable is not made supportable by a second conversation,
and the second route of \S1 exists so that a no is a redirection rather than a
dead end.

If the answers are yes but question 4 cannot be answered because a platform is
unavailable, the study is not blocked: the configuration is a site-agreement item
and a different site may hold a different platform. What that outcome does block
is the regulatory submission, which cannot name a configuration that has not been
identified, so it is recorded as an open item rather than as a failure.

\subsection{The single approval step}

\actionbox{One decision is open on this day}{Approve naming a target week for a
feasibility meeting, and approve asking two institutions and two foundations for
related things on the same day. Naming a week is the whole of it: a request with
no proposed date becomes a thread, and a request with a proposed week becomes a
calendar entry or a clear no. Both are more useful than a thread. Every letter in
\appfile{funding/auto-fund/04Sep26/emails} carries a bracket for that week and
must not be sent with the bracket unresolved.}
